%% ════════════════════════════════════════════════════════════════════════════
\section{Security and Privacy}
\label{sec:security}
%% ════════════════════════════════════════════════════════════════════════════

\subsection{Security Architecture Overview}
\label{subsec:sec_overview}

The threat model assumes deployment on a university intranet accessible
only to faculty and staff on the campus LAN.  The primary assets are
student biometric encodings and attendance records.  Defence is
structured across four layers: network (LAN-only exposure), application
(JWT role-based access), data (bcrypt, BYTEA embeddings, no raw images),
and transport (nginx as the single ingress point).

\subsection{Authentication and Authorisation}
\label{subsec:auth}

Professor accounts are authenticated via JSON Web Tokens (HS256
algorithm) with a configurable expiry (default: 480~minutes).  Passwords
are stored as bcrypt hashes; no plaintext credential ever reaches the
database.  Two roles are supported: \texttt{professor} and
\texttt{admin}.  Role-based access control is enforced at the FastAPI
dependency level (\texttt{get\_current\_professor},
\texttt{require\_admin}):

\begin{itemize}
  \item Professors see only their own courses, sessions, and
    students in all list endpoints.
  \item Admin-only endpoints (\texttt{POST /api/at-risk/recompute},
    \texttt{POST /api/forecasting/recompute}) return HTTP~403 for
    professor-role tokens.
\end{itemize}

The \texttt{REQUIRE\_AUTH} environment variable is set to \texttt{false}
in development, allowing unauthenticated requests to simplify local
testing.  This flag must be set to \texttt{true} in any deployment
accessible beyond the immediate developer machine.

\subsection{Biometric Data Handling}
\label{subsec:biometrics}

Face encodings are stored exclusively as 128-dimensional float32 numpy
arrays serialised to PostgreSQL BYTEA.  The enrollment workflow
explicitly discards the source image after computing the embedding: only
the averaged embedding vector is persisted.  No identifiable facial
image is ever written to disk or transmitted over the network.

In stub mode (\texttt{FACE\_RECOGNITION\_ENABLED=false}), a
zeroed 128-d placeholder is stored instead; the DeepFace library is not
loaded and no camera frames are processed.  This mode is used in all
Docker-based deployments, ensuring that the development environment
handles no real biometric data.

\subsection{Data Privacy Considerations}
\label{subsec:privacy}

All system components run on-premises (Raspberry Pi or Docker host).
No student data---attendance records, sensor readings, or face
encodings---is transmitted to any external service.  Ollama processes
LLM prompts locally; no student information leaves the network boundary.

The at-risk prompt engine enforces content constraints that prevent the
LLM from generating statements about a student's health, personal life,
family situation, or psychological state.  Output is limited to
attendance statistics and academic context.  The generated text is
stored in \texttt{at\_risk\_explanations.summary\_explanation} and is
visible only to professors of the student's enrolled courses and to
administrators.

Moodle synchronisation occurs over the campus LAN only.  The Moodle
token is stored as an environment variable and is never logged or
exposed through any API endpoint.

\subsection{Transport and Network Security}
\label{subsec:transport}

The current deployment uses HTTP and unencrypted WebSocket (WS) on the
university LAN.  This is acceptable for an intranet system where network
traffic is not exposed to the public internet.  The nginx reverse proxy
acts as the sole ingress point; the FastAPI backend port~8000 is not
published to the host network.

A clear path to HTTPS/WSS exists: replacing the nginx configuration
with an SSL-terminating virtual host (using a LAN certificate from a
university CA) requires no backend changes, as FastAPI operates
identically behind SSL termination.  This is documented as a
production-hardening requirement.

MQTT communication between the ESP32 and Mosquitto operates at QoS~0
over the local WiFi network.  No MQTT authentication is configured
(\texttt{allow\_anonymous true} in \texttt{mosquitto.conf}); adding
username/password or TLS client certificates is identified as future
work.

\subsection{Session and Token Security}
\label{subsec:token_security}

JWT tokens are stored in browser \texttt{localStorage}.  The security
tradeoff is acknowledged: \texttt{localStorage} is accessible to
JavaScript and is not protected by the \texttt{HttpOnly} cookie flag,
making it theoretically vulnerable to XSS attacks.  On a closed
university intranet with a controlled codebase, this risk is accepted
in exchange for simpler client-side implementation.  Token expiry
(480~minutes) is enforced server-side.  No refresh token mechanism is
implemented; after expiry the professor must re-authenticate.

\subsection{Operational Security}
\label{subsec:ops_security}

The \texttt{SECRET\_KEY} environment variable ships with a placeholder
value (\texttt{changeme-use-a-long-random-string-in-production}) that
must be replaced before any deployment.  The seed script creates
professor accounts with non-trivial passwords that are hashed with
bcrypt before insertion.  Docker named volumes isolate persistent data;
volume backup is recommended as part of a production runbook.
